% Pursuing Stacks (cotes 134-2 to 134-8): archive cote 134-8, site page 6 (PDF page 7), typescript page 560.
% typescript. Draft compared: ps7.tex:243-309.
% AI-transcribed by Claude Opus 5.5 (claude-opus-5-5, adaptive thinking, effort high), Message Batches API, batch ps4 (msgbatch_012Aoc1NuskATPsyBeKSYK7J), 2026-11-12; not yet checked by eye.
\scanpage{134-8/6}
\mspage{560}

T\hand{h}e other "extra feature"  is about  the relationship of $W^{ab}_{A}$  with
the notion of homotopism.  In case $A = \hand{\Delta}$ ,  a map in $A^{\wedge}_{ab}$  is a weak equi-%
valence (or equivalently,  an abelian weak equivalence)  iff it is a
homotopism when forgetting the abelian structures  -  this follows from
the well known fact that semisimplicial groups are Kan complexes.  There
is of course also a notion \xout{of} of homotopism in the stronger abelian sen-%
se -  a particular case of the notion of a homotopism between semisimpli-%
cial objects in an arbitrary category (here in (Ab)).  The Kan-Dold-Pupp\unsure{e}
theory implies that if F and G in $\hand{\Delta}^{\wedge}_{ab}$  have as values projective $\underline{Z}$-modu-%
les,  then \over{A}{a} map  $F \hand{\longrightarrow} G$ in $\hand{\Delta}^{\wedge}_{ab}$  is a weak equivalence iff it is an
"abelian" homotopism;  and likewise,  two maps  $u,v\colon F \hand{\rightrightarrows} G$  are equal
in $\mathrm{Hotab}_{\hand{\Delta}}$  iff they are homotopic (in the strict,\uline{abelian} sense of the
word).  T\hand{h}e corresponding statements still make a sense and are true,
when replacing $\underline{Z}$ by any ground ring k ,  working in $\hand{\Delta}^{\wedge}_{k}$  rather than
$\hand{\Delta}^{\wedge}_{\underline{Z}} = \hand{\Delta}^{\wedge}_{ab}$   -  and more generally still,  when working in
\[ \hand{\Delta}^{\wedge}_{\hand{\mathcal{M}}} = \underline{\mathrm{Hom}}(\hand{\Delta}^{o}, \hand{\mathcal{M}}) \quad , \]
where $\hand{\mathcal{M}}$ is any abelian category.  Now \deleted{\xout{for}} \inserted[typed]{replacing \hand{$\Delta$} by} an arbitrary object A in
(Cat),  these statements still make a sense, \inserted[typed]{it would seem,} provided we got on $\hand{A}^{\wedge}$ a
suitable "homotopy structure" ,  more specifically,  a suitable "homoto-%
py interval structure" (in the sense of section 51).  In section 97
(p. 355) we reviewed the three standard homotopy interval structures
which may be introduced on a category $A^{\wedge}$,  and their relationships -
the main impression remaining was that,  in case A is a \uline{contractor}
(cf. section 95),  those three structures coincide,  and may be defined
also in terms of a \uline{contractibility structure} on $A^{\wedge}$ (section 51), \inserted{the latter} giving
rise (via the general construction of section 79) to the usual notion
of weak equivalence $W_{A}$ in $A^{\wedge}$.  Thus we may hope that the feature just
mentioned  for $\hand{\Delta}^{\wedge}$ may be valid too for $A^{\wedge}$,  whenever A is a contractor.

In case A is only a test category,  \xout{the reasonable} then the most reaso-%
nable homotopy interval structure on $A^{\wedge}$ I would think,  which possibly
may still yiled the desired "extra feature",  is the one defined in \xout{\unsure{term}}
\margin{$\lfloor$}terms of the set $W_{A}$  of weak equivalences in $A^{\wedge}$ as in section 54 (namely
$h_{W_{A}}$ ) .

\margin{$\lceil$}\hand{135.}  This handful of questions (some of which we met with before) \xout{in}
\xout{a slightly different light} is mainly a way of coming into touch again
with \xout{\unsure{some}} the abelianization story,  which has been becoming somewhat
remote during the previous two months.  I am not sure I am going to

\margin{6}
